% cases/floats.tex with every float [H], float's 'here, definitely', and a float's caption set as the pipeline sets it: what the strip of that case is to be.
\documentclass{article}
\usepackage{float}
\usepackage{algorithm}
\newfloat{program}{tbp}{lop}
\floatname{program}{Program}
% a float's caption: a paragraph, small and centred
\makeatletter
\long\def\@makecaption#1#2{\vskip\abovecaptionskip{\small\centering#1: #2\par}\vskip\belowcaptionskip}
\makeatother
\pagestyle{empty}
\begin{document}
A paragraph before a figure of its own page, long enough to break onto a second
line of the column, quickly.
\begin{figure}[H]
  \centering
  \rule{3cm}{1cm}
  \caption{A figure of its own page, set where it is written.}
\end{figure}
A paragraph after the figure, long enough to break onto a second line of the
column, before a table placed here.
\begin{table}[H]
  \centering
  \begin{tabular}{ll}
    left & right \\
    gyp & quickly
  \end{tabular}
  \caption{A table placed here, with a caption long enough to need two lines
    of the column, so that it is set as a paragraph.}
\end{table}
A paragraph before a float of a type the document defines, long enough to
break onto a second line of the column.
\begin{program}[H]
  \centering
  \rule{4cm}{0.5cm}
  \caption{A program, a float of the document's own type.}
\end{program}
A paragraph before an algorithm, long enough to break onto a second line of the
column, quickly.
\begin{algorithm}[H]
  \caption{An algorithm, ruled above and below its caption}
  First step of the algorithm, long enough to break onto a second line of the
  column, quickly.

  Second step.
\end{algorithm}
A paragraph before a starred figure, long enough to break onto a second line of
the column, ending gyp.
\begin{figure}[H]
  \centering
  \rule{5cm}{1cm}
  \caption{A starred figure, as wide as the text.}
\end{figure}
A last paragraph, after all the floats, long enough to break onto a second line
of the column.
\end{document}
